\chapter{Справочные величины для Zr--H}

Значения --- ориентиры для оценок; для расчётов конкретного сплава берите его собственные измерения.

\begin{center}\footnotesize
\renewcommand{\arraystretch}{1.25}
\begin{tabularx}{\linewidth}{>{\raggedright\arraybackslash}p{4.3cm}>{\raggedright\arraybackslash}p{5.6cm}>{\raggedright\arraybackslash}X}
\toprule
величина & значение & источник, комментарий \\
\midrule
решётка $\alpha$-Zr (ГПУ) & $a = 0.323$~нм, $c = 0.515$~нм, $c/a = 1.593$ & комнатная температура \\
превращение $\alpha\to\beta$ в чистом Zr & 863\degC & \\
$\delta$-гидрид & ГЦК, ZrH$_{1.5\ldots1.7}$, $a\approx0.478$~нм & \\
ориентационное соотношение & $(0001)_\alpha\parallel\{111\}_\delta$ & плоскость залегания около $\{10\bar17\}$ (Westlake, 1968) \\
несоответствие $\delta$-гидрида & $\eps_a\approx0.046$, $\eps_c\approx0.072$, объём $\approx17\,\%$ & Carpenter (1973) \\
водород в гидриде & $\approx18\,000$~ppm (ZrH$_{1.66}$) & \\
объём гидрида на моль H & $\approx1.0\cdot10^{-5}$~м$^3$ & из плотности ZrH$_{1.66}$ \\
TSSD, Э635 & $\lg C[\text{ppm}] = 5.56 - 2240/T$ & ДСК, Плясов и др. (2023); $Q\approx43$~кДж/моль \\
TSSP, Э635 & $\lg C[\text{ppm}] = 5.08 - 1680/T$ & то же; $Q\approx32$~кДж/моль; зависит от режима \\
эвтектоид $\beta\to\alpha+\delta$ & $\approx550$\degC, $\approx6$~ат.\,\% H в $\alpha$ & Zuzek и др. (1990) \\
диффузия H в $\alpha$-Zr & $D = 7.9\cdot10^{-7}e^{-44.4~\text{кДж}/RT}$~м$^2$/с & Kearns (1972) \\
парциальный мольный объём H & $\bar V_{\mathrm H}\approx1.7$~см$^3$/моль & \\
упругость Zr (300--400\degC) & $E\approx80$--90~ГПа, $\nu\approx0.34$, $\mu\approx30$--33~ГПа & изотропное приближение \\
предел текучести оболочки (300--400\degC) & 200--500~МПа & от состояния (рекр./СОН) \\
энергия межфазной границы гидрида & $\sim0.1$--0.3~Дж/м$^2$ & плохо известна; параметр моделей \\
сдвиг выпадения под нагрузкой & $\approx0.08$~К/МПа & Zr--2.5Nb (Vizca\'{\i}no и др., 2014) \\
факторы Кернса оболочек & $f_r\approx0.5$--0.7, $f_\theta\approx0.25$--0.35 & \\
\bottomrule
\end{tabularx}
\end{center}

\section*{Перевод единиц}
\begin{itemize}
\item ppm (масс.) $\to$ ат.\,\%: $x = \dfrac{C/1.008}{C/1.008 + 10^6/91.22}$; 100~ppm $\approx0.9$~ат.\,\%.
\item $RT$ при 300\degC: 4.76~кДж/моль; $k_BT = 0.049$~эВ. 1~эВ на атом = 96.5~кДж/моль.
\item 1~МПа = 1~МДж/м$^3$; на моль водорода в гидриде: 1~МДж/м$^3$ $\approx10$~Дж/моль~H.
\end{itemize}
